% Übungsklausur Diskrete Mathematik (Schema der Klausur 29.07.2026, neue Zahlen).
% Alle Zahlen mit den Werkzeugen in skripte/ berechnet und geprüft
% (rsa, inklusion_exklusion, euklid, crt, mod11_code).
\providecommand{\teil}[1]{\par\medskip\textbf{#1)}\ }

\section{Übungsklausur}
\subsection{Aufgaben}

\begin{aufgabe}[Aufgabe 1]
Alice will die Nachricht $w = 920$ verschlüsseln und als verschlüsselte Nachricht $c$ (Codewort)
an Bob senden. Sie verabreden als Modul die Zahl $m = 47 \cdot 59 = 2773$ zu wählen und den
öffentlichen Schlüssel $e = 157$ zu verwenden.
\begin{enumerate}[label=\alph*)]
  \item Schreiben Sie $e$ im Binärsystem.
  \item Mit welcher Formel verschlüsselt Alice das Wort $w$? Berechnen Sie das Codewort $c$.
  \item Bob besitzt den privaten Schlüssel $d$. Wie (Formel) entschlüsselt Bob das Codewort $c$?
  \item Berechnen Sie $\varphi(m)$.
  \item Oscar erhält Kenntnis von $m$ und $e$. Welche Formeln muss er nutzen, um den privaten
        Schlüssel $d$ zu berechnen? Berechnen Sie $d$ auf zwei Wegen
        -- mit dem Euklidischen Algorithmus bzw.\ der Matrix $Q$ und mit dem Satz von Euler:
        \[ d \equiv e^{\varphi(\varphi(m)) - 1} \bmod \varphi(m) . \]
        Machen Sie die Probe $c^d \equiv w \pmod m$.
\end{enumerate}
\end{aufgabe}

\begin{aufgabe}[Aufgabe 2]
Wie viele positive ganze Zahlen $\leq 840$ sind weder durch $3$, $5$ oder $7$ teilbar?
Verwenden Sie die Symbole $A_i$, $\alpha_i$ aus der Vorlesung. (Inklusion -- Exklusion)
\end{aufgabe}

\begin{aufgabe}[Aufgabe 3]
Bestimmen Sie mit Hilfe des Euklidischen Algorithmus
\begin{enumerate}[label=\alph*)]
  \item den größten gemeinsamen Teiler von $m = 29$ und $n = 23$,
  \item eine ganzzahlige Lösung der Gleichung $29x - 23y = 1$ (EA rückwärts und Matrix $Q$),
  \item sämtliche ganzzahligen Lösungen dieser Gleichung.
\end{enumerate}
\end{aufgabe}

\begin{aufgabe}[Aufgabe 4]
Gegeben ist das System simultaner Kongruenzen (Chinesischer Restsatz)
\begin{align*}
  z &\equiv 4 \pmod{29}\\
  z &\equiv 7 \pmod{23}
\end{align*}
\begin{enumerate}[label=\alph*)]
  \item Bestimmen Sie eine ganze Zahl $z$ mit Hilfe des Ergebnisses aus Aufgabe 3.
  \item Bestimmen Sie sämtliche Lösungen dieses Systems, insbesondere die kleinste positive Lösung.
\end{enumerate}
\end{aufgabe}

\begin{aufgabe}[Aufgabe 5]
Die Parameter eines Modulo-11-Codes sind $[10, 8, 3]_{11}$.
Vorgelegt ist das Klartextwort $a = 2\,0\,3\,1\,0\,0\,0\,4$.
\begin{enumerate}[label=\alph*)]
  \item Wie lautet die Kontrollmatrix $H$?
  \item Berechnen Sie die Kontrollziffern $c_9$, $c_{10}$ und damit das zugehörige Codewort $c$.
\end{enumerate}
Empfangen wurden die Wörter (die Ziffer $10$ steht für eine Stelle)
\begin{align*}
  r_1 &= 2\,0\,3\,1\,0\,0\,3\,4\,2\,10\\
  r_2 &= 2\,0\,3\,1\,1\,1\,0\,4\,2\,10
\end{align*}
\begin{enumerate}[label=\alph*), start=3]
  \item Berechnen Sie die Syndrome $S(r_1)$, $S(r_2)$.
  \item Welches der beiden Wörter $r_1$, $r_2$ ist nicht dekodierbar und warum nicht?
  \item Dekodieren Sie das andere Wort und geben Sie das Codewort $c$ an.
  \item Zeigen Sie, dass $c$ tatsächlich ein Codewort ist.
\end{enumerate}
\end{aufgabe}

\clearpage
% ===========================================================================
\section{Lösungen der Übungsklausur}

\begin{center}
\begin{tabular}{@{}llp{0.52\textwidth}@{}}
\toprule
Aufgabe & Thema & Ergebnis\\
\midrule
1 & RSA & $e = 10011101_2$, $c = 192$, $\varphi(m) = 2668$, $d = 17$\\
2 & Inklusion -- Exklusion & $384$\\
3 & Euklidischer Algorithmus & $\ggT(29,23) = 1$, $x = 4 + 23t$, $y = 5 + 29t$\\
4 & Chinesischer Restsatz & $z \equiv 352 \pmod{667}$, kleinste positive Lösung $352$\\
5 & Modulo-11-Code & $c = 2\,0\,3\,1\,0\,0\,0\,4\,2\,10$; $r_1 \to c$ (Stelle 7); $r_2$ nicht dekodierbar\\
\bottomrule
\end{tabular}
\end{center}

% ---------------------------------------------------------------------------
\subsection*{Aufgabe 1}
Gegeben: $w = 920$, $m = 47 \cdot 59 = 2773$, $e = 157$.

\teil{a} \textbf{$e$ im Binärsystem.}
\[
157 = 128 + 16 + 8 + 4 + 1 = 2^7 + 2^4 + 2^3 + 2^2 + 2^0
\quad\Rightarrow\quad e = 10011101_2 .
\]

\teil{b} \textbf{Verschlüsseln.} Formel (Alice): $c \equiv w^e \pmod m$, also $c \equiv 920^{157} \pmod{2773}$.

Schnelles Potenzieren: $920^{157} = 920^{1} \cdot 920^{4} \cdot 920^{8} \cdot 920^{16} \cdot 920^{128}$.

Quadrattabelle (mod 2773, jede Zeile ist das Quadrat der Zeile davor):
\[
\begin{array}{r@{\;\equiv\;}l@{\qquad}l}
920^{1}   & 920 & \leftarrow\\
920^{2}   & 920^2 = 846400 = 305 \cdot 2773 + 635 \equiv 635 & \\
920^{4}   & 635^2 = 403225 = 145 \cdot 2773 + 1140 \equiv 1140 & \leftarrow\\
920^{8}   & 1140^2 = 1299600 = 468 \cdot 2773 + 1836 \equiv 1836 & \leftarrow\\
920^{16}  & 1836^2 = 3370896 = 1215 \cdot 2773 + 1701 \equiv 1701 & \leftarrow\\
920^{32}  & 1701^2 = 2893401 = 1043 \cdot 2773 + 1162 \equiv 1162 & \\
920^{64}  & 1162^2 = 1350244 = 486 \cdot 2773 + 2566 \equiv 2566 & \\
920^{128} & 2566^2 = 6584356 = 2374 \cdot 2773 + 1254 \equiv 1254 & \leftarrow
\end{array}
\]
Produkt Schritt für Schritt (mod 2773):
\begin{align*}
920 \cdot 1140 &= 1048800 = 378 \cdot 2773 + 606 \equiv 606\\
606 \cdot 1836 &= 1112616 = 401 \cdot 2773 + 643 \equiv 643\\
643 \cdot 1701 &= 1093743 = 394 \cdot 2773 + 1181 \equiv 1181\\
1181 \cdot 1254 &= 1480974 = 534 \cdot 2773 + 192 \equiv 192
\end{align*}
\begin{ergebnis} $c = 920^{157} \bmod 2773 = 192$ \end{ergebnis}

\teil{c} \textbf{Entschlüsseln.} Formel (Bob): $w \equiv c^d \pmod m$.
Das funktioniert, weil $e \cdot d \equiv 1 \pmod{\varphi(m)}$ und nach Euler
$c^d \equiv (w^e)^d = w^{ed} \equiv w \pmod m$.

\teil{d} \textbf{$\varphi(m)$.} $47$ und $59$ sind Primzahlen (Probedivision: $47$ ist durch $2, 3, 5$ nicht teilbar und $7^2 > 47$;
$59$ ist durch $2, 3, 5, 7$ nicht teilbar und $8^2 > 59$). Also
\[
\varphi(2773) = \varphi(47)\,\varphi(59) = (47 - 1)(59 - 1) = 46 \cdot 58 = 2668 .
\]
(Produktformel: $2773 \cdot \tfrac{46}{47} \cdot \tfrac{58}{59} = 2668$.)
\begin{ergebnis} $\varphi(m) = 2668$ \end{ergebnis}

\teil{e} \textbf{Privater Schlüssel $d$.}
Oscar zerlegt $m = 47 \cdot 59$, berechnet $\varphi(m) = 2668$ und löst
\[
e \cdot d \equiv 1 \pmod{\varphi(m)}, \quad\text{also}\quad 157 \cdot d \equiv 1 \pmod{2668}.
\]

\emph{Weg 1: Euklidischer Algorithmus (Lemma von Bezout).}
\[
\begin{array}{r@{\;=\;}l}
2668 & 16 \cdot 157 + 156\\
157 & 1 \cdot 156 + 1\\
156 & 156 \cdot 1 + 0
\end{array}
\qquad\Rightarrow\quad \ggT(2668, 157) = 1 .
\]
EA rückwärts:
\[
1 = 157 - 156 = 157 - (2668 - 16 \cdot 157) = 17 \cdot 157 - 1 \cdot 2668 .
\]
Matrix $Q = \prod \begin{pmatrix} q_i & 1\\ 1 & 0\end{pmatrix}$ mit $q = 16, 1, 156$:
\[
\begin{pmatrix}16&1\\1&0\end{pmatrix}\begin{pmatrix}1&1\\1&0\end{pmatrix} = \begin{pmatrix}17&16\\1&1\end{pmatrix},\qquad
\begin{pmatrix}17&16\\1&1\end{pmatrix}\begin{pmatrix}156&1\\1&0\end{pmatrix} = \begin{pmatrix}2668&17\\157&1\end{pmatrix} = Q .
\]
$\det Q = 2668 \cdot 1 - 157 \cdot 17 = 2668 - 2669 = -1$, also $157 \cdot 17 - 2668 \cdot 1 = 1$, d.\,h.
$157 \cdot 17 \equiv 1 \pmod{2668}$ und $d = 17$.

\emph{Weg 2: Satz von Euler.} Aus $e^{\varphi(\varphi(m))} \equiv 1 \pmod{\varphi(m)}$ folgt
\[
d \equiv e^{\varphi(\varphi(m)) - 1} \pmod{\varphi(m)} .
\]
Zerlegung $2668 = 2 \cdot 1334 = 2^2 \cdot 667 = 2^2 \cdot 23 \cdot 29$, also
\[
\varphi(2668) = \varphi(2^2)\,\varphi(23)\,\varphi(29) = 2 \cdot 22 \cdot 28 = 1232,
\qquad d \equiv 157^{1231} \pmod{2668}.
\]
$1231 = 1024 + 128 + 64 + 8 + 4 + 2 + 1 = 10011001111_2$.
Quadrattabelle (mod 2668):
\[
\begin{array}{r@{\;\equiv\;}l@{\qquad}l}
157^{1}    & 157 & \leftarrow\\
157^{2}    & 157^2 = 24649 = 9 \cdot 2668 + 637 \equiv 637 & \leftarrow\\
157^{4}    & 637^2 = 405769 = 152 \cdot 2668 + 233 \equiv 233 & \leftarrow\\
157^{8}    & 233^2 = 54289 = 20 \cdot 2668 + 929 \equiv 929 & \leftarrow\\
157^{16}   & 929^2 = 863041 = 323 \cdot 2668 + 1277 \equiv 1277 & \\
157^{32}   & 1277^2 = 1630729 = 611 \cdot 2668 + 581 \equiv 581 & \\
157^{64}   & 581^2 = 337561 = 126 \cdot 2668 + 1393 \equiv 1393 & \leftarrow\\
157^{128}  & 1393^2 = 1940449 = 727 \cdot 2668 + 813 \equiv 813 & \leftarrow\\
157^{256}  & 813^2 = 660969 = 247 \cdot 2668 + 1973 \equiv 1973 & \\
157^{512}  & 1973^2 = 3892729 = 1459 \cdot 2668 + 117 \equiv 117 & \\
157^{1024} & 117^2 = 13689 = 5 \cdot 2668 + 349 \equiv 349 & \leftarrow
\end{array}
\]
Produkt Schritt für Schritt (mod 2668):
\begin{align*}
157 \cdot 637 &= 100009 = 37 \cdot 2668 + 1293 \equiv 1293\\
1293 \cdot 233 &= 301269 = 112 \cdot 2668 + 2453 \equiv 2453\\
2453 \cdot 929 &= 2278837 = 854 \cdot 2668 + 365 \equiv 365\\
365 \cdot 1393 &= 508445 = 190 \cdot 2668 + 1525 \equiv 1525\\
1525 \cdot 813 &= 1239825 = 464 \cdot 2668 + 1873 \equiv 1873\\
1873 \cdot 349 &= 653677 = 245 \cdot 2668 + 17 \equiv 17
\end{align*}
Beide Wege liefern $d = 17$. Kontrolle: $157 \cdot 17 = 2669 = 1 \cdot 2668 + 1 \equiv 1 \pmod{2668}$ \checkmark
\begin{ergebnis} $d = 17$ \end{ergebnis}

\textbf{Probe} $c^d \equiv w$: $192^{17} = 192^{16} \cdot 192^{1}$, Quadrattabelle (mod 2773):
\[
\begin{array}{r@{\;\equiv\;}l@{\qquad}l}
192^{1}  & 192 & \leftarrow\\
192^{2}  & 192^2 = 36864 = 13 \cdot 2773 + 815 \equiv 815 & \\
192^{4}  & 815^2 = 664225 = 239 \cdot 2773 + 1478 \equiv 1478 & \\
192^{8}  & 1478^2 = 2184484 = 787 \cdot 2773 + 2133 \equiv 2133 & \\
192^{16} & 2133^2 = 4549689 = 1640 \cdot 2773 + 1969 \equiv 1969 & \leftarrow
\end{array}
\]
\[
192 \cdot 1969 = 378048 = 136 \cdot 2773 + 920 \equiv 920 = w \quad\checkmark
\]

% ---------------------------------------------------------------------------
\subsection*{Aufgabe 2}
$\Omega = \{1, \dots, 840\}$, $|\Omega| = 840$. Sei $A_i = \{k \in \Omega \mid i \text{ teilt } k\}$ für $i = 3, 5, 7$.
Gesucht ist
\[
|\Omega \setminus (A_3 \cup A_5 \cup A_7)| = |\Omega| - \alpha_1 + \alpha_2 - \alpha_3 .
\]
\begin{align*}
\alpha_1 &= |A_3| + |A_5| + |A_7| = \tfrac{840}{3} + \tfrac{840}{5} + \tfrac{840}{7} = 280 + 168 + 120 = 568,\\
\alpha_2 &= |A_3 \cap A_5| + |A_3 \cap A_7| + |A_5 \cap A_7| = \tfrac{840}{15} + \tfrac{840}{21} + \tfrac{840}{35} = 56 + 40 + 24 = 120,\\
\alpha_3 &= |A_3 \cap A_5 \cap A_7| = \tfrac{840}{105} = 8 .
\end{align*}
(Schnittmengen über das kgV: $\kgV(3,5) = 15$, $\kgV(3,7) = 21$, $\kgV(5,7) = 35$, $\kgV(3,5,7) = 105$.)
\[
840 - 568 + 120 - 8 = 384 .
\]
Kontrolle: $840 \cdot (1 - \tfrac13)(1 - \tfrac15)(1 - \tfrac17) = 840 \cdot \tfrac23 \cdot \tfrac45 \cdot \tfrac67 = 384$.
\begin{ergebnis} $384$ Zahlen $\le 840$ sind durch keine der Zahlen $3, 5, 7$ teilbar. \end{ergebnis}

% ---------------------------------------------------------------------------
\subsection*{Aufgabe 3}
\teil{a}
\[
\begin{array}{r@{\;=\;}l}
29 & 1 \cdot 23 + 6\\
23 & 3 \cdot 6 + 5\\
6 & 1 \cdot 5 + 1\\
5 & 5 \cdot 1 + 0
\end{array}
\qquad\Rightarrow\quad \ggT(29, 23) = 1 .
\]

\teil{b} EA rückwärts (von unten nach oben einsetzen):
\begin{align*}
1 &= 6 - 5\\
  &= 6 - (23 - 3 \cdot 6) = 4 \cdot 6 - 23\\
  &= 4 \cdot (29 - 23) - 23 = 4 \cdot 29 - 5 \cdot 23 .
\end{align*}
Also $29 \cdot 4 - 23 \cdot 5 = 116 - 115 = 1$.

Matrix $Q$ mit den Quotienten $q = 1, 3, 1, 5$:
\begin{align*}
\begin{pmatrix}1&1\\1&0\end{pmatrix}\begin{pmatrix}3&1\\1&0\end{pmatrix} &= \begin{pmatrix}4&1\\3&1\end{pmatrix},\qquad
\begin{pmatrix}4&1\\3&1\end{pmatrix}\begin{pmatrix}1&1\\1&0\end{pmatrix} = \begin{pmatrix}5&4\\4&3\end{pmatrix},\\
\begin{pmatrix}5&4\\4&3\end{pmatrix}\begin{pmatrix}5&1\\1&0\end{pmatrix} &= \begin{pmatrix}29&5\\23&4\end{pmatrix} = Q .
\end{align*}
$\det Q = 29 \cdot 4 - 23 \cdot 5 = 1$ ($= (-1)^{n+1}$ mit $n = 3$), also dieselbe Lösung $x = 4$, $y = 5$.

\teil{c} Homogene Gleichung $29x - 23y = 0$ hat die Lösung $(23, 29)$. Alle Lösungen:
\[
x = 4 + 23t, \qquad y = 5 + 29t \qquad (t \in \Z).
\]
Probe: $29 \cdot 4 - 23 \cdot 5 = 116 - 115 = 1$ \checkmark
\begin{ergebnis} $\ggT(29,23) = 1$; eine Lösung $x = 4$, $y = 5$; alle Lösungen $x = 4 + 23t$, $y = 5 + 29t$ ($t \in \Z$). \end{ergebnis}

% ---------------------------------------------------------------------------
\subsection*{Aufgabe 4}
\teil{a} Ansatz: $z = 29x + 4 = 23y + 7$, also
\[
29x - 23y = 7 - 4 = 3 .
\]
Aus Aufgabe 3: $29 \cdot 4 - 23 \cdot 5 = 1$. Multiplikation mit $3$:
\[
29 \cdot 12 - 23 \cdot 15 = 3, \qquad x = 12,\ y = 15 .
\]
\[
z = 29 \cdot 12 + 4 = 352 = 23 \cdot 15 + 7 .
\]

\teil{b} Da $\ggT(29, 23) = 1$, sind alle Lösungen $z \equiv 352 \pmod{29 \cdot 23 = 667}$:
\[
z = 352 + k \cdot 667 \quad (k \in \Z).
\]
Kleinste positive Lösung: $z = 352$ (denn $352 - 667 < 0$).

Probe: $352 = 12 \cdot 29 + 4 \equiv 4 \pmod{29}$ \checkmark, \quad $352 = 15 \cdot 23 + 7 \equiv 7 \pmod{23}$ \checkmark
\begin{ergebnis} $z \equiv 352 \pmod{667}$, also $z = 352 + 667k$ ($k \in \Z$); kleinste positive Lösung $z = 352$. \end{ergebnis}

% ---------------------------------------------------------------------------
\subsection*{Aufgabe 5}
\teil{a} Kontrollmatrix (Stellen $1, \dots, 10$):
\[
H = \begin{pmatrix}
1 & 1 & 1 & 1 & 1 & 1 & 1 & 1 & 1 & 1\\
1 & 2 & 3 & 4 & 5 & 6 & 7 & 8 & 9 & 10
\end{pmatrix}
\]
Ein Wort $c$ ist Codewort genau dann, wenn $H c^T = 0$, also
$\sum_{i=1}^{10} c_i \equiv 0$ und $\sum_{i=1}^{10} i \cdot c_i \equiv 0 \pmod{11}$.

\teil{b} Kontrollziffern für $a = 2\,0\,3\,1\,0\,0\,0\,4$ ($c_i = a_i$ für $i = 1, \dots, 8$):
\begin{align*}
c_9 &\equiv \sum_{i=1}^{8} (1+i)\, c_i = 2 \cdot 2 + 4 \cdot 3 + 5 \cdot 1 + 9 \cdot 4\\ &= 57 = 5 \cdot 11 + 2 \equiv 2 \pmod{11},\\
c_{10} &\equiv \sum_{i=1}^{8} (9-i)\, c_i = 8 \cdot 2 + 6 \cdot 3 + 5 \cdot 1 + 1 \cdot 4\\ &= 43 = 3 \cdot 11 + 10 \equiv 10 \pmod{11}.
\end{align*}
\begin{ergebnis} $c_9 = 2$, $c_{10} = 10$, $c = 2\,0\,3\,1\,0\,0\,0\,4\,2\,10$ \end{ergebnis}

\teil{c} Syndrome $S(r) = H r^T = (s_1, s_2)^T$.

$r_1 = 2\,0\,3\,1\,0\,0\,3\,4\,2\,10$:
\begin{align*}
s_1 &= 2 + 3 + 1 + 3 + 4 + 2 + 10 = 25 \equiv 3,\\
s_2 &= 2 + 3 \cdot 3 + 4 \cdot 1 + 7 \cdot 3 + 8 \cdot 4 + 9 \cdot 2 + 10 \cdot 10\\
    &= 2 + 9 + 4 + 21 + 32 + 18 + 100\\
    &= 186 = 16 \cdot 11 + 10 \equiv 10
\quad\Rightarrow\quad S(r_1) = (3, 10)^T .
\end{align*}
$r_2 = 2\,0\,3\,1\,1\,1\,0\,4\,2\,10$:
\begin{align*}
s_1 &= 2 + 3 + 1 + 1 + 1 + 4 + 2 + 10 = 24 \equiv 2,\\
s_2 &= 2 + 3 \cdot 3 + 4 \cdot 1 + 5 \cdot 1 + 6 \cdot 1 + 8 \cdot 4 + 9 \cdot 2 + 10 \cdot 10\\
    &= 2 + 9 + 4 + 5 + 6 + 32 + 18 + 100\\
    &= 176 = 16 \cdot 11 \equiv 0
\quad\Rightarrow\quad S(r_2) = (2, 0)^T .
\end{align*}

\teil{d} Bei genau einem Fehler der Größe $e_x$ an der Stelle $x$ gilt $S = e_x \cdot (1, x)^T$,
also $e_x = s_1$ und $x = s_2 / s_1$.

Für $r_2$: $x = 0 / 2 = 0 \cdot 2^{-1} = 0 \cdot 6 \equiv 0$. Die Stelle $0$ gibt es nicht (Stellen $1, \dots, 10$).
Das Syndrom passt also zu keinem Einzelfehler: $r_2$ enthält mindestens zwei Fehler
(tatsächlich an den Stellen $5$ und $6$, je $+1$: $s_1 = 1 + 1 = 2$, $s_2 = 5 + 6 = 11 \equiv 0$).
Der Code hat $d = 3$ und korrigiert nur $1$ Fehler.
\begin{ergebnis} $r_2$ ist nicht dekodierbar (mindestens 2 Fehler). \end{ergebnis}
\begin{achtung}
Zwei Fehler sind nicht immer erkennbar: Wären die Stellen $9$ und $10$ je um $1$ verfälscht
($r = 2\,0\,3\,1\,0\,0\,0\,4\,3\,0$), so wäre $S = (2, 19)^T \equiv (2, 8)^T$ und $x = 8 \cdot 2^{-1} = 8 \cdot 6 = 48 \equiv 4$.
Das sieht aus wie ein Einzelfehler an Stelle $4$ und wird \emph{falsch} zu $2\,0\,3\,10\,0\,0\,0\,4\,3\,0$ decodiert.
\end{achtung}

\teil{e} Für $r_1$: $e_x = s_1 = 3$ und $x = s_2 / s_1 = 10 \cdot 3^{-1} = 10 \cdot 4 = 40 \equiv 7 \pmod{11}$
(denn $3 \cdot 4 = 12 \equiv 1$). Fehler an Stelle $7$ mit Größe $3$:
\[
c_7 = r_7 - 3 = 3 - 3 = 0 .
\]
\begin{ergebnis} $c = 2\,0\,3\,1\,0\,0\,0\,4\,2\,10$ \end{ergebnis}

\teil{f} Probe $H c^T = 0$:
\begin{align*}
\text{Zeile 1:}\quad & 2 + 3 + 1 + 4 + 2 + 10 = 22 = 2 \cdot 11 \equiv 0 \pmod{11},\\
\text{Zeile 2:}\quad & 2 + 9 + 4 + 32 + 18 + 100 = 165 = 15 \cdot 11 \equiv 0 \pmod{11}.
\end{align*}
Also ist $c$ ein Codewort. \checkmark

\medskip
{\small\textit{Hinweis:} Andere Formen von $H$ (z.\,B. systematisch $H = (A \mid E)$) beschreiben denselben Code.
Die Syndromwerte sehen dann anders aus, das Ergebnis (Fehlerstelle, Codewort) ist dasselbe.}
